
\documentclass{report}

\input{../../preamble}
\input{../../macros}
\input{../../letterfonts}

\title{IEOR 173: Problem Set 3}
\author{Ryan Lin}
\date{}

\begin{document}
\maketitle
\section{Textbook Questions}

\qs{5}{An urn contains 3 white, 6 red, 5 black balls. Six of these balls are randomly selected from the urn. Let $ X $ and $ Y $ denote respectively the number of white and black balls selected. compute the conditional pmf of $ X $ given that $ Y=3 $  Also compute $ E[X \mid Y=1] $
\begin{align*}
	p_{X|Y}(x\mid y=3) &= P(X=x | Y=3) \\
			   &= \frac{P(X=x, Y=3)}{P(Y=3)} \\
			   &= \frac{\binom{3}{k}\binom{6}{3-k}}{\binom{9}{3}} \\
\end{align*}


\begin{align*}
	E[X|Y=1] &= \sum_{k=0}^{5} k \cdot P(X=k \mid Y=1) \\
	 &= \sum_{k=0}^{3} k \cdot \frac{\binom{3}{k} \binom{6}{6-k}}{\binom{9}{5}} \\
	 &= 1 \cdot \frac{3 \cdot 15}{126}+ 2 \cdot \frac{3 \cdot 20}{126} + 3 \cdot \frac{1 \cdot 15}{126}\\
	 &= \frac{210}{126}
\end{align*}
}

\qs{8}{
An unbiased die is successively rolled. Let $ X $ and $ Y $ denote the number of rolls necessary to obtain a six and a five respectively. 
\vspace{5mm}

a.) Find $ E[X] $ 

$ X $ and $ Y $ are random variables with distribution $ geometric(\frac{1}{6}) $ 

\[
E[X] = 6
\]


\vspace{5mm}
b.) Find $ E[X \mid Y=1] $ 

\[
E[X \mid Y=1] = 1 + E[X] = 7
\]
\vspace{5mm}
c.) Find $ E[X \mid Y=5] $ 
\begin{align*}
	E[X \mid Y=5] &= 1 + P(X < y \mid Y=y)E[X \mid X<y, Y=y] + P(X > y \mid Y=y)E[X] \\
	&= 1 + \sum_{k=1}^{4} k \cdot  P(X = k \mid Y=y) + (\frac{4}{5})^{y-1}(5 + E[X]) \\
	&= 1 + \sum_{k=1}^{4} k \cdot  \left( \frac{4}{5} \right)^{k-1} \cdot  \frac{1}{5}+ 11\left( \frac{4}{5} \right) ^{4}  \\
	&\approx 5.8192 \\
\end{align*}
}

\qs{11}{
The joint density of $ X $ and $ Y $  is given by 
\[
	f(x,y)  = \frac{(y^{2}-x^{2})}{8}e^{-y}, \quad -y<x<y, 0<y<\infty  
\]

Show that $ E[X \mid Y=y] = 0 $ 

\begin{align*}
	f_Y(y) &= \int_{-y}^{y} \frac{(y^{2}-x^{2})}{8}e^{-y} dx \\
	&= \frac{e^{-y}}{8} \int_{-y}^{y} y^{2} - x^{2} dx  \\
	&= \frac{e^{-y}}{8} \left[ xy^{2} - \frac{1}{3}x^{3} \right]_{-y}^{y} \\
	&= \frac{e^{-y}}{8}\cdot \left( \frac{4}{3}y^{3} \right)  \\
	&= \frac{1}{6}y^{3}e^{-y} \\
\end{align*}

\begin{align*}
	f_{X\mid Y}(x,y) &= \frac{f(x,y)}{f_Y(y)} \\
	&= \frac{\frac{1}{8}(y^{2} -x^{2})e^{-y}}{\frac{1}{6}y^{3}e^{-y}} \\
	&= \frac{6}{8} \cdot \frac{y^{2} -x^{2}}{y^{3}} \\
	&= \frac{3}{4} \cdot \frac{y^{2}- x^{2}}{y^{3}} \\
\end{align*}

Since $ f_{X\mid Y}(x,y) $ is an even function when $ y $ is fixed, and the bounds of $ x $ are symmetrical, we know that the expected value has to be 0.  However, this still does not mean that $ X $ and $ Y $ are independent. 

}
\newpage
\qs{21}{Consider the example of a miner trapped in a mine. Let $ N $ denote the total number of doors selected before the miner reaches safety. Also let $ T_i $ denote the travel time corresponding to the $ i $th choice, $ i\ge 1 $. Again, let $ X $ denote the time the miner reaches safety. 

a.) Give an identity that relate $ X $ to $ N $ and the $ T_i $ 

\[
X  = \sum_{i=1}^{N} T_i
\]

b.) What is $ E[N] $ 

Let $ Y $ indicate what door the miner initially chose. 
\[
	E[N] = E[N \mid Y=1]P(Y=1) + E[N \mid Y=2]P(Y=2) + E[N \mid Y=3]P(Y=3)
\]
\begin{align*}
	E[N \mid Y=1] &= 1 \\
	E[N \mid Y=2] &= 1 + E[N] \\
	E[N \mid Y=3] &= 1 + E[N] \\
\end{align*}

\begin{align*}
	E[N] &= \frac{1}{3} (1 + 1 + E[N] + 1 + E[N] \\
	&= \frac{1}{3} (3 + 2E[N]) \\
	&= 1 + \frac{2}{3} E[N] \\
	\frac{1}{3} E[N] &= 1 \\
	E[N] &= 3 \\
\end{align*}

c.) What is $ E[T_i] $ 

\[
E[T_i] = 2P(Y=1) + 3P(Y=2) + 5P(Y=3) = \frac{1}{3} (2+3+5) = \frac{10}{3} \\
\]

d.) What is $ E[\sum_{i=1}^{N} T_i \mid N=n] $ 

\begin{align*}
	E[\sum_{i=1}^{N}T_i \mid N=n] &= \sum_{i=1}^{N} E[T_i \mid N = n] \\
	&= \sum_{i=1}^{N-1} E[T_i \mid N=n] + E[T_n \mid N=n] \\
	&= \sum_{i=1}^{N-1} E[T_i \mid \text{Door 2 or Door 3}] + E[T_i \mid \text{Door 1}] \\
	&= (N-1) \cdot \frac{1}{2} (3+5) + 2 \\
	&= 4N - 4 + 2 \\
	&= 4N - 2 \\
\end{align*}


e.) Using the preceding, what $ E[X] $ 

\[
E[X] = E[E[X \mid N]] = E[E[\sum_{i=1}^{N}T_i \mid N=n]] = E[4N-2] = 4E[N] -2 = 10
\]
}

\qs{26}{
You have 2 opponents with whom you alternate play. Whenever you play $ A $, you with with probability $ p_a $. When you play $ B $ you win with probability $ p_b $, where $ p_b > p_a $. Your objective is to min expected number of games to win 2 in a row. Should you start with A or B?

Let $ N_i $  be the number of games until you get 2 in a row if you start with player $ i $ .

\[
E[N_a] = 1 + p_a(p_b(1) + (1-p_b)(1+E[N_a]) + (1-p_a)E[N_b] = 1 + p_ap_b + p_a(1-p_b)(1+ E[N_a]) + (1-p_a)E[N_b]
\]
\[
	E[N_b] = 1 + p_b(p_a(1) + (1-p_a)(1+E[N_b]) + (1-p_b)E[N_a] = 1 + p_a p_b + p_b(1-p_a)(1+E[N_b]) + (1-p_b)E[N_a]
\]
After subtracting both of these equations and simplifying: 

\[
E[N_a] -E[N_b]  = \frac{p_a - p_b}{(2-p_a -p_b +p_ap_b)}
\]

Therefore, since $ p_a >p_b $, then the RHS is negative. Therefore, that implies that $ E[N_b] > E[N_a]$, which means you should start with $ A $  
}

\qs{40}{
A prisoner is trapped in a cell containing three doors. The first door leads to a
tunnel that returns him to his cell after two days of travel. The second leads to a
tunnel that returns him to his cell after three days of travel. The third door leads
immediately to freedom.

\vspace{5mm}
(a) Assuming that the prisoner will always select doors 1, 2, and 3 with prob
abilities 0.5, 0.3, 0.2, what is the expected number of days until he reaches
freedom?

Let $ X $ be the number of days to reach freedom. 

\begin{align*}
	E[X] &= 0.5(2 + E[X]) + 0.3(3 + E[X]) + 0.2(0) \\
	&= 1 + 0.5E[X] + .9 + .3E[X] \\
	&= 1.9 + 0.8E[X] \\
	0.2E[X] &= 1.9 \\
	E[X] &= 9.5
\end{align*}


\vspace{5mm}
(b) Assuming that the prisoner is always equally likely to choose among those
doors that he has not used, what is the expected number of days until he
reaches freedom? (In this version, for instance, if the prisoner initially tries
door 1, then when he returns to the cell, he will now select only from doors
2 and 3.)

\[
E[X] = E[X \mid \text{door 1 chosen first}] + E[X \mid \text{door 2 chosen first}] + E[X \mid \text{door 3 chosen first}]
\]
\[
E[X \mid \text{door 1 chosen first}] = 2 + \frac{1}{2} \cdot 3 = 3.5
\]
\[
E[X \mid \text{door 2 chosen first}]= 3+ \frac{1}{2} \cdot 2 = 4
\]

\[
E[X \mid \text{door 3 chosen first}] = 0
\]

\[
E[X] = \frac{1}{3} (3.5 + 4) = 2.5 
\]
\vspace{5mm}
(c) For parts (a) and (b) find the variance of the number of days until the prisoner
reaches freedom.
\[
	\Var(X) = E[X^2] - E[X]^{2}
\]
\begin{align*}
	E[X^{2}] &= .5(E[(2+X)^{2}] + 0.3(E[(3+X)^{2}]) \\ 
	&= \frac{1}{2}(4 +4E[X] + E[X^{2}] + .3(9+6E[X] + E[X^{2}] \\
	&= 2 + 2E[X] + .5E[X^{2}] + 2.7 + 1.8E[X] + 0.3E[X^{2}] \\
	&= 4.7 + 3.8E[X] + .8E[X^{2}] \\
	&= 204 \\
\end{align*}

\[
\Var(X) = 204 - 9.5^{2} = 113.75
\]

\begin{align*}
	\Var(X) &= \sum_{k=0}^{5} k^{2}P(X=x) \\
	&= 2 \cdot \frac{1}{6} + 3 \cdot \frac{1}{6} + 5 \cdot \frac{1}{6} + 5 \cdot \frac{1}{6} \\
	&= 4.25 \\
\end{align*} 
}

\newpage

\section{Fun Questions}
\qs{Financial Engineering}{
Suppose that Joe invests equal amoutns in $ n $ different independently performing stocks. Let $ Z = \sum_{i=1}^{n}\frac{X}{n} $ be his avg return, where $ X_i $ is the return for stock i. Let $ Y_i  = X_i - Z$ be the deviation in return for stock i from the average, and assume all stock have the same standard deviation. 

a.) Show that $\Cov(Y, Z) = 0 $ 


\[
\Cov(Z, Y_i) = \Cov(Z, X_i - Z) = \Cov(Z, X_i) - \Cov(Z, Z) = \Cov(Z, X_i) - \Var(Z)
\]
\[
\Cov(Z,X_i) = \frac{1}{n} \sum_{n=1}^{n}\Cov(X_j, X_i)
\]
Because Stocks are independent, then  this covariance is 0 for $ i\neq j $ and $ \Var(X) $ for $ i=j $   

\[
\Cov(Z, Y_i) = \frac{1}{n}\Var(X) - \Var(Z) = \frac{1}{n}\Var(X) - \frac{1}{n^{2}}\sum_{}^{}\Var(X) = \frac{1}{n}\Var(x) - \frac{n\Var(X)}{n^{2}} = 0
\]

\vspace{5mm}
b.) 

\[
\Var(Y_i) = \Var(X_i - Z) = \Var(X_i) + \Var(Z) - 2\Cov(X_i, Z) = \Var(X_i) + \frac{\Var(X_i)}{n} - 2 \frac{\Var(X_i)}{n} = \Var(X) (\frac{n-1}{n})
\]
\[
\Var(X) = \Var(Y) \frac{n}{n-1}
\]

\[
\Var(Z) = \frac{\Var(x)}{n} = \Var(Y) \cdot \frac{1}{n-1}
\]
}


\qs{Operations Management}{
If the defect rate of a machine is 10\%, where defective widgets are
produced independently, which do you think will be larger: the probability of at least 2 defective
widgets in a box of 10 widgets or the probability of at least 20 defective widgets in a box of 100
widgets? Now check your intuition by doing the math. (You can use Excel or an on-line binomial
calculator.)

I think the probability of at least 2 widgets in a box of 10 widgets, because when the number of trials is lower, the variance is higher. 


For the box of 10: 

\[
P(X=0) = \binom{10}{0}0.1^{0}0.9^{10} = 0.3487, P(X=1) = \binom{10}{1}0.1^{1}0.9^{9} = 0.3874
\]
\[
P(X\ge_2) = 1-P(X=0) - P(X=1) = 0.2639
\]

Using a calculator for 100: 

\[
P(X\ge 20) = 0.00198
\]
}

\qs{Telecommunications}{
Consider a telecommunications channel that can transmit 1 packet in
each time slot. The numbers of packet arrivals during each slot are i.i.d., where the number of
arrivals is 0 with probability .5, 1 with probability .4 and more than 1 with probability .1. Think of
arrivals as occurring just at the end of a slot, so if a slot starts with 1 packet, and there is 1 arrival,
the first packet will be transmitted and the next slot will start with 1 packet. The buffer can only
hold 2 packets, so if a slot starts with 2 packets and 2 packets arrive, one packet will be transmitted,
one of the arriving packets will be lost, and the next slot will start with 2 packets. Find the mean
length of a busy period (the time from when the channel goes from idle to busy until it goes back to
idle again).


Let $ N $  be the length of the busy period and $ N_i $ where $ i $ be the number of packets in storage during the busy period. Let $ A $ be the number of packets arriving

\begin{align*}
	E[N_1] &= E[N_1 \mid A = 0]P(A=0) + E[N_1\mid A=1]P(A=1) + E[N_1 \mid A>1]P(A>1)\\
	&= 1(0.5) + (1+E[N_1])0.4 + (1+ E[N_2])0.1 \\
	&= 0.5 + 0.4 + 0.4E[N_1] + 0.3 + 0.1E[N_1] \\
	&= 1.2 +0.5E[N_1] \\
	&= 1.2  \\
	&= 2.4 \\
\end{align*}


\[
E[N_2] = E[N_2 \mid A=0]0.5 + E[N_2 \mid A=1]0.4 +E[N_2 \mid A>1]0.1
\]
\[
E[X_2] = (1+E[N_1]).5 + (1+E[N_2]).4 + (1+E[N_2]).1 
\]
\[
= .5 + .5E[N_1] + .5 + .5E[N_2] \rightarrow E[X_2] =2 + E[N_1] = 4.4
\]

\[
E[N]= \frac{P(A=1)}{P(A>0)}E[N_1] + \frac{P(A>1)}{P(A>0)}E[N_2] = 0.8(2.4) + 0.2(4.4)
\]
}

\end{document}
